\chapter{Ketones (Blood)}

\section*{What Is This Test?}
Blood ketone measurement quantifies beta-hydroxybutyrate or total ketones in serum or plasma. Beta-hydroxybutyrate is the predominant ketone body in diabetic ketoacidosis. The body produces ketones (acetoacetate, beta-hydroxybutyrate, acetone) through hepatic beta-oxidation of fatty acids. They serve as alternative fuel when glucose is scarce. The test is used to diagnose and manage diabetic ketoacidosis (DKA) and alcoholic ketoacidosis. It also monitors ketogenic diet therapy for epilepsy. Blood testing supersedes urine ketone detection. It measures beta-hydroxybutyrate directly, the principal ketone in DKA. It correlates better with clinical status \cite{dhatariya2020diabetic}.

\section*{Alternative Names}
\begin{itemize}
  \item Beta-Hydroxybutyrate
  \item BHB
  \item Blood Ketones
  \item Serum Ketones
  \item Plasma Ketones
  \item Acetoacetate
  \item Beta-Hydroxybutyric Acid
\end{itemize}
\section*{Principle}
An enzymatic method runs on automated chemistry analyzers. It uses beta-hydroxybutyrate dehydrogenase. The reaction is: beta-hydroxybutyrate + NAD+ $\leftrightarrow$ acetoacetate + NADH + H+. The rate of NADH formation is measured at 340 nm. It is proportional to beta-hydroxybutyrate concentration. Point-of-care (POC) meters include Abbott Precision Xtra and KetoScan. They use similar enzyme-based test strips. They provide results within 10 seconds from a fingerstick \cite{arora2011point}.

\section*{History}
Ketonuria was first detected by Gerhardt's ferric chloride test (1865). Modern enzymatic quantification of beta-hydroxybutyrate was introduced in the 1970s. Testing shifted from urine acetoacetate to blood beta-hydroxybutyrate. This reflected the finding that in DKA beta-hydroxybutyrate outnumbers acetoacetate 3:1. During recovery from DKA, acetoacetate shifts back to BHB. So urine ketones paradoxically increase.

\section*{Why This Test Is Ordered}
\begin{itemize}
  \item It is used for emergency evaluation of hyperglycemia with metabolic acidosis (suspected DKA).
  \item It helps diagnose alcoholic ketoacidosis in alcohol use disorder with vomiting and acidosis.
  \item It monitors DKA treatment. The aim is to halve beta-hydroxybutyrate every 2 hours with insulin.
  \item It evaluates hypoglycemia with ketosis. Causes include inborn errors of metabolism and starvation.
  \item It monitors ketogenic diet therapy. The target blood BHB is 1.5--3.0 mmol/L for seizure control \cite{neal2008ketogenic}.
  \item It helps differentiate DKA from hyperosmolar hyperglycemic state (HHS).
\end{itemize}

\section*{Primary vs Secondary Test}
It is the primary test in suspected DKA. It is also primary in managing ketoacidosis.

\section*{How This Test Is Performed}
Venous blood goes into a serum or heparin tube. POC testing uses a fingerstick. Ketones are volatile. They also convert reversibly between BHB and acetoacetate. So the sample is processed quickly. Lab results take 30--60 minutes. POC results take seconds.

\section*{What Preparation Is Needed}
None is needed in emergency situations. For ketogenic diet monitoring, take a fasting morning blood sample. This keeps results consistent.

\section*{Contraindications and Precautions}
\begin{itemize}
  \item Urine ketone (nitroprusside) tests measure acetoacetate, not beta-hydroxybutyrate. They may paradoxically increase during DKA recovery despite clinical improvement. Capillary blood ketone meters and strips should be reasonably calibrated. Falsely low readings occur with high acetone or prolonged storage. POC equipment is less reliable than laboratory measurement in severe DKA.
\end{itemize}
\section*{Risks}
Risks are limited to fingerstick or venipuncture.

\section*{What Happens After the Test}
In suspected DKA, repeat beta-hydroxybutyrate every 2--4 hours. Do this with electrolytes, glucose, and ABG/venous blood gas. The treatment goal is BHB $<$ 1.5 mmol/L. Anion gap closure is also a goal.

\section*{Understanding Results}

\subsection*{Below Normal}
Low or normal ketones with ketosis signs can occur if acetoacetate predominates. This happens in alcoholic ketoacidosis and in recovery from DKA. Use clinical judgment.

\subsection*{Above Normal}
\begin{itemize}
  \item \textbf{Diabetic ketoacidosis:} BHB is $>$ 3.0 mmol/L. It results from insulin deficiency and counterregulatory hormone excess \cite{kitabchi2009hyperglycemic}.
  \item \textbf{Alcoholic ketoacidosis:} BHB is 2--10 mmol/L. It results from dietary depletion of glycogen and alcohol metabolism.
  \item \textbf{Starvation ketosis:} BHB is 0.5--3.0 mmol/L. It is a physiologic adaptation and is usually benign.
  \item \textbf{Ketogenic diet therapy:} BHB is targeted at 1.5--3.0 mmol/L for seizure control.
  \item \textbf{Inborn errors of metabolism:} Examples include beta-ketothiolase deficiency and carnitine palmitoyltransferase deficiencies.
  \item \textbf{SGLT2 inhibitors (euglycemic DKA):} Glucose is normal or mildly high with high BHB.
\end{itemize}

\section*{Normal Results}
\begin{itemize}
  \item \textbf{Serum beta-hydroxybutyrate (fasting, normal):} 0.02--0.27 mmol/L.
  \item \textbf{After overnight fast:} up to 0.6 mmol/L.
  \item \textbf{Ketotic threshold:} 0.5--1.0 mmol/L (mild ketosis).
  \item \textbf{Ketoacidosis:} $>$ 3.0 mmol/L is probable. $>$ 4.8 mmol/L is confirmatory.
  \item \textbf{Anion gap:} $>$ 12 mEq/L with high BHB and metabolic acidosis indicates DKA.
\end{itemize}

\section*{Follow-up Tests}
\begin{itemize}
  \item \textbf{Venous blood gas (VBG):} pH $<$ 7.30, bicarbonate $<$ 15 mEq/L confirms acidosis.
  \item \textbf{Blood glucose:} $>$ 250 mg/dL is typical in DKA. Euglycemic DKA is possible with SGLT2 inhibitors.
  \item \textbf{Serum electrolytes and anion gap:} the anion gap is $>$ 12 mEq/L.
  \item \textbf{Serum beta-hydroxybutyrate (serial):} measure every 2 hours during DKA management.
  \item \textbf{Serum osmolality:} it differentiates DKA ($<$ 320 mOsm/kg) from HHS ($>$ 320 mOsm/kg).
  \item \textbf{Lactate, alcohol level, salicylate, acetaminophen:} these test for other causes of high-anion-gap acidosis.
\end{itemize}

\section*{References}
\bibliographystyle{unsrtnat}
\bibliography{references}

\clearpage
\section*{Regulatory Status and Authorization Date}
This chapter describes a test category, not a specific commercial product. FDA, CDSCO, EU IVDR, and MHRA approval or registration dates therefore cannot be assigned without the assay manufacturer, product name, instrument, and intended use. For laboratory-developed tests, an individual agency approval date may not exist. See the Preface for the interpretation of regulatory status.

\clearpage
